\documentclass[11pt,letterpaper]{article}

\usepackage[utf8]{inputenc}
\usepackage[spanish]{babel}
\usepackage[margin=2.5cm]{geometry}
\usepackage{amsmath,amssymb}
\usepackage{enumitem}
\usepackage{titlesec}
\usepackage{array}
\usepackage{hyperref}
\usepackage{xcolor}

\definecolor{unicaucaverde}{RGB}{0,90,60}
\hypersetup{colorlinks=true,linkcolor=unicaucaverde,urlcolor=unicaucaverde}

\titleformat{\section}{\normalfont\Large\bfseries\color{unicaucaverde}}{\thesection}{1em}{}
\titleformat{\subsection}{\normalfont\large\bfseries}{\thesubsection}{1em}{}

\newcommand{\vv}[1]{\mathbf{#1}}

\setlength{\parindent}{0pt}
\setlength{\parskip}{6pt}

\title{\textcolor{unicaucaverde}{\Large Guía de Aprendizaje — Semana 3}\\[4pt]
  \large Vectores y matrices: definiciones y operaciones básicas}
\author{Álgebra Lineal para Ingeniería --- Universidad del Cauca}
\date{2026}

\begin{document}
\maketitle

\begin{center}
\begin{tabular}{|>{\bfseries}p{4cm}|p{10cm}|}
\hline
Semana & 3 de 16 \\ \hline
Unidad & 2 --- Vectores y matrices (Grossman, cap. 2) \\ \hline
Subtemas & 2.1 (vectores y matrices), 2.2 (productos vectorial y matricial) \\
\hline
Horas de clase & 4 \\ \hline
Resultado de aprendizaje & RA2 (parte 1 de 3; ver \texttt{guia1/guia.tex}) \\ \hline
\end{tabular}
\end{center}

\section*{Relevancia para ingeniería}

Vectores y matrices son el lenguaje natural para representar varias cantidades relacionadas a la
vez: fuerzas y momentos en un plano, señales de varios canales, tablas de datos, transformaciones
geométricas (rotación, escalado) en gráficos por computador, o insumos y recursos en un proceso de
producción. El producto matricial en particular condensa en una sola operación todo un conjunto de
combinaciones lineales; es la operación base de casi todo el software numérico de ingeniería y del
aprendizaje automático (una red neuronal es, en esencia, una cadena de productos matriciales). Esta
semana conecta además con la Unidad 1: todo sistema de ecuaciones lineales $A\vv{x}=\vv{b}$ es, en
el fondo, un producto matricial.

\section{Objetivo de la semana}

Representar cantidades mediante vectores y matrices, y calcular sus operaciones básicas (suma,
producto por escalar, producto punto, producto matricial), incluida la representación matricial
$A\vv{x}=\vv{b}$ de un sistema de ecuaciones lineales.

\section{Resultados de aprendizaje de la semana}

\begin{itemize}[leftmargin=1.2cm]
  \item[\textbf{RA2.1.}] Definir vector y matriz, usar la notación de doble subíndice $a_{ij}$,
        reconocer matrices especiales (fila, columna, cuadrada, nula, identidad, diagonal), y
        efectuar suma, resta y multiplicación por escalar de matrices.
  \item[\textbf{RA2.2.}] Calcular el producto punto de dos vectores y el producto de matrices
        (incluido matriz por vector), verificar cuándo un producto está definido según las
        dimensiones de los factores, y expresar un sistema de ecuaciones lineales en la forma
        matricial $A\vv{x}=\vv{b}$.
\end{itemize}

\section{Contenidos}

\begin{itemize}[leftmargin=1.2cm]
  \item \textbf{2.1.} Vectores y matrices: definiciones básicas, notación, igualdad, matrices
        especiales, suma, resta y multiplicación por escalar.
  \item \textbf{2.2.} Productos vectorial y matricial: producto punto, producto matriz-vector,
        producto de matrices y sus propiedades.
\end{itemize}

\section{Plan de la sesión (4 horas)}

\begin{enumerate}[leftmargin=1.2cm]
  \item[\textbf{Hora 1.}] Definición de vector (fila/columna) y de matriz $m\times n$; notación
        $a_{ij}$; igualdad de matrices; matrices especiales (fila, columna, cuadrada, nula,
        identidad, diagonal) (teoría + ejemplos).
  \item[\textbf{Hora 2.}] Suma, resta y multiplicación por escalar de matrices; propiedades
        algebraicas. Taller: operar con matrices dadas y verificar propiedades.
  \item[\textbf{Hora 3.}] Producto punto de dos vectores; producto matriz-vector $A\vv{x}$ como
        combinación lineal de columnas; forma matricial $A\vv{x}=\vv{b}$ de un sistema de
        ecuaciones (conecta con la Unidad 1).
  \item[\textbf{Hora 4.}] Producto de matrices $AB$ (condición de compatibilidad, regla fila por
        columna) y sus propiedades, incluida la no conmutatividad en general. Taller de
        refuerzo con los seis ejercicios de \texttt{notas.tex}.
\end{enumerate}

\section{Actividades}

\begin{itemize}[leftmargin=1.2cm]
  \item \textbf{En clase}: talleres de las horas 2 y 4 (ver arriba).
  \item \textbf{Trabajo autónomo}: lista de ejercicios de \texttt{notas.tex} §1--2, de entrega
        antes de la Semana 4.
\end{itemize}

\section{Material de apoyo}

\begin{itemize}[leftmargin=1.2cm]
  \item \texttt{notas.tex} --- desarrollo teórico completo de 2.1 y 2.2, con ejemplos y ejercicios.
  \item \texttt{diapositivas.tex} --- síntesis visual de la sesión.
\end{itemize}

\section{Evaluación de la semana}

Control formativo corto al inicio de la Semana 4 (10 minutos): dadas dos o tres matrices, calcular
una suma, un producto por escalar y un producto matricial, y escribir un sistema pequeño en forma
$A\vv{x}=\vv{b}$, para verificar RA2.1 y RA2.2 antes de avanzar a 2.3--2.5. Contribuye al Quiz de
la Unidad 2 (ver \texttt{guia1/guia.tex} § Evaluación).

\section{Bibliografía de la semana}

\begin{enumerate}[leftmargin=1.2cm]
  \item Grossman, S. I., Flores Godoy, J. J. \emph{Álgebra lineal}. 7a/8a ed., McGraw-Hill,
        cap. 2, §2.1--2.2.
  \item Lay, D. C., Lay, S. R., McDonald, J. J. \emph{Álgebra lineal y sus aplicaciones}. Pearson.
\end{enumerate}

\end{document}
